\chapter{总结与展望}
\echapter{Conclusion and Future Work}

\section{全文总结}
\esection{Conclusion}

随着数字经济的蓬勃发展，数据要素的跨机构融合与共享已成为释放数据核心价值、驱动各行业创新的关键引擎。然而，在日益严格的数据合规监管背景下，“数据孤岛”与“隐私保护”之间的矛盾愈发凸显。多方隐私集合运算（Multi-party Private Set Operations, MPSO）作为安全多方计算（Secure Multi-party Computation, MPC）中极具应用前景的重要分支，为打破数据壁垒、实现多方数据的“可用不可见”提供了基础支撑。然而，由于多方场景相较于两方场景存在着极高的组合复杂性，长期以来，MPSO 领域的研究深陷于安全性弱、效率低下、功能单一局限以及缺乏统一理论框架等一系列核心困局。

针对上述挑战，本文立足于大规模复杂数据协作的真实需求，开展了系统性且不断深入的研究。从突破特定功能的性能与安全瓶颈，到构建支持任意复合集合运算的统一框架，再到实现从通用理论到专精工程的性能转化，本文形成了一套涵盖底层密码学原语创新、协议架构设计与系统工程验证的完整多方隐私集合运算技术体系。全文的核心研究脉络与主要贡献总结如下：

首先，本文着眼于应用价值巨大但研究严重滞后的多方隐私集合求并（Multi-party Private Set Union, MPSU）方向，彻底解决了该领域长期存在的“安全性与实用性难以兼顾”的难题。针对现有基于对称密钥技术的最高效 MPSU 协议 \cite{LG-ASIACRYPT-2023} 在安全性上严重依赖“非共谋假设”（Non-collusion Assumption）的致命缺陷，本文提出并构造了批量秘密分享隐私成员测试（batch secret-shared private membership test, batch ssPMT）和多方秘密分享随机不经意传输（multi-party secret-shared ROT, mss-ROT）两个轻量级密码学原语。基于此为组件，本文设计了首个在标准半诚实模型下基于对称密钥的高效 MPSU 协议（SK-MPSU），在成功消除共谋攻击风险的同时，实现了相较于现有最先进方案 \cite{LG-ASIACRYPT-2023} 成倍的性能跃升。

其次，本文直击基于秘密分享技术路线的 MPSU 框架所固有的超线性复杂度瓶颈 —— 由于最终获得结果集合的领导者（Leader）需要对每个元素执行多方份额重构，其计算与通信复杂度必然随参与方数量呈二次增长。为了实现真正适用于大规模多方协作的线性复杂度目标，本文巧妙地将 batch ssPMT 原语与多密钥重随机化公钥加密（Multi-Key Rerandomizable Public-Key Encryption, MKR-PKE）技术相融合，提出了首个同时实现线性计算复杂度和线性通信复杂度的 MPSU 协议（PK-MPSU），大幅降低了带宽受限环境下的通信开销。随后，本文在真实的局域网与广域网环境下对 SK-MPSU 与 PK-MPSU 进行了系统的基准测试与多维对比，量化了两种协议在不同网络拓扑和数据规模下的性能边界，为现实世界中带宽受限或算力受限场景下的协议部署提供了科学的选型依据。

然后，为了打破现有 MPSO 研究中呈现的严重“功能孤岛”现象，满足金融风控、医疗流调等复杂业务中对交、并、差任意组合的复合集合运算需求，本文构建了首个兼具理论完备性与实际可用效率的通用 MPSO 统一框架。在表示层，本文创新性地引入了规范集合谓词公式（Canonical Set Predicate Formula, CSPF）作为任意集合公式的统一代数表示；在原语层，提出了具有极强抽象与组合能力的“谓词零分享”（Predicative Zero-Sharing）密码学原语，并给出其在集合运算下的特化实例化——“成员零分享”（Membership Zero-Sharing）。通过自底向上的模块化构造范式，该统一框架彻底摆脱了以往通用方案对昂贵通用 MPC 计算电路的依赖，不仅高效实现了任意集合公式计算的通用 MPSO 功能，还将框架扩展至多方隐私集合求势（MPSO-card）和基于电路的通用 MPSO（Circuit-MPSO）这两个更加复杂但通用的计算功能，为此前缺乏高效可用方案的复杂集合统计场景提供了理论支撑和技术支持。同时，作为对上述基于对称密钥技术路线的补充，本文还提出了“谓词零加密（Predicative Zero-Encryption）”原语，并构建了基于公钥体制的通用 MPSO 框架，进一步探索了通用 MPSO 框架在渐近复杂度上的理论边界，共同构筑了完整的 MPSO 理论生态。

最后，本文成功弥合了“协议通用性”与“专用高效性”之间的鸿沟，打通了从理论框架到工程落地的最后一公里。本文深入剖析了特定集合公式（如交集、并集）内在结构特性，对各个典型 MPSO 具体功能进行了深度的协议级优化。优化后的各实例化协议在理论层面取得了重大突破：多方隐私集合求交（Multi-party Private Set Intersection, MPSI）、多方隐私交集求势（MPSI-card）及多方隐私交集求势与和（MPSI-card-sum）协议首次在标准半诚实模型下实现了与明文计算同量级的最优渐近复杂度。在实践层面，本文通过系统级的全面实现与对比实验证明，在该统一框架下优化的各实例化协议，其实际在线运行效率已达到甚至超越了现有各细分领域内的专用 SOTA 协议。这标志着本文所构建的 MPSO 技术体系不仅在理论框架上实现了从从“零散孤立”到“统一通用”的跨越，在工程实践中同样确立了领先优势，为破解数字经济时代复杂多变的数据安全协同难题提供了坚实、普适且极致高效的密码学底层基础设施。

\section{未来工作展望}
\esection{Future Work}

尽管本文在标准半诚实模型下构建了安全、高效且功能完备的多方隐私集合运算框架，并解决了一系列理论与工程上的开放性问题，但在极其复杂的开放网络和存在巨大利益冲突的真实商业合作中，参与方有时具有更强的攻击动机。他们可能不仅会试图从协议执行的中间信息中推断隐私，还可能恶意偏离协议规定——例如篡改输入、拒绝响应或注入恶意消息以破坏结果的正确性并窃取数据。在这些高敏感场景（如金融风控对抗、医疗数据确权）中，半诚实模型提供的安全保证往往是不够的。因此，研究能够抵御恶意敌手任意偏离行为的\textbf{恶意安全（Malicious Security）}的隐私集合运算协议，具有极其重要的理论意义与实际应用价值，这将是本文未来工作的核心聚焦方向。

当前，恶意安全的隐私集合运算（PSO）研究现状呈现出明显的两极分化：在两方 PSO 领域，恶意安全的隐私集合求交（PSI）协议已经相当成熟，涌现出了一系列高效的方案 \cite{RT-CCS-2021, RS-EUROCRYPT-2021, RR-CCS-2022, RindalR17}，覆盖了从小集合到大规模集合的各类应用场景。目前最先进的两方 PSI 协议已经能够做到只需付出极小的额外计算与通信开销，即可将协议的安全性从半诚实模型提升至恶意安全模型。然而，相较之下，恶意安全的两方隐私交集计算（PCSI） \cite{MiaoP0SY20} 和两方隐私集合求并（PSU） \cite{PuGT26} 协议由于构造难度极大，相关研究极少，且仅存的协议效率极低，根本无法适用于现实的工业级应用。特别值得注意的是在 PSU 方向，Hazay 和 Nissim 曾深刻指出 \cite{HN-JOC-2012}：在单向输出（仅一方获得结果）的恶意安全设定下，接收方实际上拥有极大的权力——恶意的接收方完全可以将其输入集合替换为整个数据全集，从而在不经意间提取出发送方的所有私有数据。由于这种攻击在理想模型下是允许的，这就使得单向输出的恶意安全  PSU 协议退化为等同于发送方直接明文发送数据的毫无意义的操作。因此，要想构造有意义的恶意安全 PSU，必须研究\textbf{双向输出（Two-Sided Output）}的恶意安全机制，这进一步大幅拉高了协议设计的门槛。

在多方场景（MPSO）下，恶意安全的协议设计更加举步维艰。虽然多方隐私集合求交（MPSI）方向已有突破——例如 Kolesnikov 等人 \cite{KMPRT-CCS-2017} 提出的增强半诚实 MPSI 框架随后被证明天然支持恶意安全 \cite{GPRTY-CRYPTO-2021}，而 Nevo 等人 \cite{NTY-CCS-2021} 进一步优化了阈值腐化假设下的恶意安全  MPSI。通过结合最先进的不经意伪随机函数（Oblivious Pseudorandom Function, OPRF）~\cite{FIPR-TCC-2005,CM-CRYPTO-2020,RS-EUROCRYPT-2021} 与不经意键值存储（Oblivious Key-Value Store, OKVS）\cite{PRTY20,RR-CCS-2022,BPSY-USENIX-2023} 组件，当前的恶意安全  MPSI 协议已经达到了极高的运行效率。同时，基于混淆布隆过滤器（Garbled Bloom Filter）的技术路线也被成功扩展至恶意安全模型 \cite{BNOP-AsiaCCS-2022}。但是，对于 MPSO 中的其他复杂功能（如 MPCSI 和 MPSU），目前学术界是完全的无人区。这主要归因于以下两大根本性阻碍：
\begin{enumerate}
\item \textbf{哈希分桶技术在恶意模型下的失效.} 本文以及目前最高效的半诚实安全 MPSO 协议广泛依赖于哈希分桶（Hashing to Bins）技术来将全局计算降维为局部计算。然而，该技术在恶意安全模型下存在致命漏洞：恶意参与方可以故意将其输入集合的元素放入错误的哈希桶中，从而人为引发匹配失败或碰撞，破坏协议正确性甚至窃取其他方输入。强行利用零知识证明来保证分桶的正确性将带来不可承受的计算灾难。
\item \textbf{模块化组合的安全性崩溃.} 半诚实安全 MPSO 协议往往由多个子协议模块化组合而成。但在恶意模型下，这种简单的组合会导致安全性失效。以本文第三章为例，利用批量不经意可编程伪随机函数（batch Oblivious Programmable Pseudorandom Function, batch OPPRF） 与 （Secret-Shared Private Equality Test, ssPEQT） 组合构造 batch ssPMT 协议在半诚实下是安全的；但在恶意设定下，即使这两个底层子协议自身具备恶意安全性，恶意敌手仍可以在协议交接处作恶（例如不将 batch OPPRF 得到的正确伪随机输出馈送给后续的 ssPEQT 协议）。要保证整个协议链条的一致性与防篡改，需要极其复杂且高效的状态验证机制，而现阶段缺乏此类紧凑的密码学转化工具。
\end{enumerate}

综上所述，未来工作计划将遵循“先两方、后多方”的递进式研究路线。首先，我们将重点攻坚恶意安全的两方 PCSI 协议以及双向输出的恶意两方 PSU 协议，致力于探索能避开昂贵非交互式零知识证明（NIZK）的新型理论框架，寻找更高效的具体构造。在掌握高效且具备抗恶意组合特性的两方组件后，我们将进一步探索突破哈希分桶安全缺陷和模块化组合陷阱的方法，最终研究从两方场景向多方复杂场景（恶意安全  MPCSI 与恶意安全  MPSU）的平滑扩展，从而为未来更加开放、恶劣的网络数据互联互通提供金融级的安全保障底座。